\subsection[Extracting R\_123 from R\_123]{Extracting $\boldsymbol{R_{123}}$ from $\boldsymbol{\widehat{R}_{123}}$}

Let us illustrate the action of the monomial part
$\tau_{\varepsilon_1, \varepsilon_2, \varepsilon_3, \varepsilon_4}$
 \eqref{taue} of $\widehat{R}_{123}$ on the canonical variables for the case
 $(\varepsilon_1, \varepsilon_2, \varepsilon_3, \varepsilon_4) = (-,-,+,+)$.
 From Example \ref{ex:1} and \eqref{Yw}--\eqref{Ypw}, we find that
$\tau_{--++}$ is translated into a transformation $\tau^{\uu\ww}_{--++}$
of the canonical variables\footnote{$\tau^{\uu\ww}_{--++}$ is naturally regarded also as
a transformation in $\mathcal{W}_3$ via exponentials.} as
\begin{align}
\tau^{\uu\ww}_{--++}\colon\
\begin{cases}
\uu_1 \mapsto \uu_1+\uu_2-\uu_3+\lambda_0,
&
\ww_1 \mapsto \ww_1+ \lambda_1,
\\
\uu_2 \mapsto \uu_3-\lambda_0,
&
\ww_2 \mapsto \ww_1+\ww_3+\lambda_3,
\\
\uu_3 \mapsto \uu_2+\lambda_0,
&
\ww_3 \mapsto -\ww_1+\ww_2+\lambda_2,
\end{cases}
\label{uwL1}
\end{align}
where $\lambda_r=\lambda_r(\mathscr{P}_1,\mathscr{P}_2,\mathscr{P}_3)$
for $r=0, 1, 2, 3$ is defined, under the condition \eqref{ae0},
by
\begin{gather}
\lambda_0 = \frac{e_2-e_3}{2},\qquad
\lambda_1 = a_2-a_3+b_2-b_3 + \lambda_0,\qquad
\lambda_2 = -a_1-b_2+b_3-\lambda_0,\nonumber\\
\lambda_3 = c_1-c_2+c_3.\label{lad}
\end{gather}

In order to realize \eqref{uwL1} as an adjoint action, we introduce
the group $N_n$ generated by
\begin{gather*}
\e^{\pm \tfrac{1}{\hbar}\uu_i\ww_j}\quad (i \neq j),\qquad
\e^{ \tfrac{a}{\hbar}\uu_i}, \qquad \e^{ \tfrac{a}{\hbar}\ww_i}  \quad (a \in \C) ,\qquad b \in \C^\times
\end{gather*}
with $i, j\in \{1,\dotsc, n\}$. The multiplication is defined by
the (generalized) Baker--Campbell--Hausdorff (BCH) formula and
\eqref{uwh}, which is well defined due to the grading by $\hbar^{-1}$.
Let $\mathfrak{S}_n$ be the symmetric group generated by the
transpositions $\rho_{ij}$ ($i, j \in \{1,\dotsc, n\}$). It acts on
$N_n$ via the adjoint action, inducing permutations of the indices of
the canonical variables. Thus one can form the semi-direct product
$N_n \rtimes \mathfrak{S}_n$, and let it act on $\mathcal{W}_n$ by the
adjoint action.

Now the monomial part $\tau^{\uu\ww}_{--++}$
is described as the adjoint action as follows:
\begin{align}
&\tau^{\uu\ww}_{--++} = \mathrm{Ad}(P_{--++}),
\label{Pbch0}\\
&P_{--++}=
\e^{\tfrac{1}{\hbar}(\uu_3-\uu_2)\ww_1}
\e^{\tfrac{\lambda_0}{\hbar}(-\ww_1-\ww_2+\ww_3)}
\e^{\tfrac{1}{\hbar}(\lambda_1\uu_1+\lambda_2\uu_2+\lambda_3\uu_3)}\rho_{23}
\in N_3 \rtimes \mathfrak{S}_3,
\label{Pbch1}
\end{align}
where $\rho_{23}$ acts trivially on the parameters $\lambda_r$.
%
Extending the indices and suppressing the sign choice of $--++$ in \eqref{Pbch0} and \eqref{Pbch1}, we introduce
\begin{align}
\tau^{\uu\ww}_{ijk} &= \mathrm{Ad}(P_{ijk}),
\label{tijk}\\
P_{ijk} &=
\e^{\tfrac{1}{\hbar}(\uu_k-\uu_j)\ww_i}
\e^{\tfrac{\lambda_0}{\hbar}(-\ww_i-\ww_j+\ww_k)}
\e^{\tfrac{1}{\hbar}(\lambda_1\uu_i+\lambda_2\uu_j+\lambda_3\uu_k)}\rho_{jk}
\in N_6 \rtimes \mathfrak{S}_6,
\label{pijk}
\end{align}
where $\lambda_r=\lambda_r(\mathscr{P}_i, \mathscr{P}_j, \mathscr{P}_k)$ is given by
\eqref{lad} by replacing the parameters as
$(\mathscr{P}_1 ,\mathscr{P}_2 ,\mathscr{P}_3) \rightarrow
(\mathscr{P}_i ,\mathscr{P}_j ,\mathscr{P}_k)$.
%
By a straightforward calculation using the BCH formula, one can prove the following.
\begin{Lemma}\label{le:tep}
$P_{ijk}$ satisfies the tetrahedron equation in $N_6 \rtimes \mathfrak{S}_6$ by itself
\begin{align}\label{pte}
P_{124}P_{135}P_{236}P_{456} =
P_{456}P_{236}P_{135}P_{124}.
\end{align}
\end{Lemma}
The fact that $P_{ijk}$ acts on the canonical variables rather than $Y$-variables
has led to the tetrahedron equation {\em without} a twist by permutations.

Define $R_{123}= R(\mathscr{P}_1,\mathscr{P}_2,\mathscr{P}_3)_{123}$ by
\begin{align}
R_{123}
={}&
\Psi_q\bigl(\e^{-d_1-c_2-b_3+\uu_1+\uu_3+\ww_1-\ww_2+\ww_3}\bigr)^{-1}
\Psi_q\bigl(\e^{-d_1-c_2-b_3-e_3+\uu_1-\uu_3+\ww_1-\ww_2+\ww_3}\bigr)^{-1}
\nonumber \\
& \times
\Psi_q\bigl(\e^{d_1+e_1+c_2+b_3+\uu_1-\uu_3-\ww_1+\ww_2-\ww_3}\bigr)\nonumber \\
& \times
\Psi_q\bigl(\e^{d_1+e_1+c_2+e_2+b_3+\uu_1+2\uu_2-\uu_3-\ww_1+\ww_2-\ww_3}\bigr)P_{123}
\label{RL0}
\\
={}&
\Psi_q\bigl(\e^{-d_1-c_2-b_3+\uu_1+\uu_3+\ww_1-\ww_2+\ww_3}\bigr)^{-1}
\Psi_q\bigl(\e^{-d_1-c_2-b_3-e_3+\uu_1-\uu_3+\ww_1-\ww_2+\ww_3}\bigr)^{-1}
\nonumber\\
& \times P_{123}
\Psi_q\bigl(\e^{-b_1-a_2-d_3-e_3+\uu_1-\uu_3+\ww_1-\ww_2+\ww_3}\bigr)\nonumber \\
& \times
\Psi_q\bigl(\e^{-b_1-a_2-d_3+\uu_1+\uu_3+\ww_1-\ww_2+\ww_3}\bigr).
\label{RL1}
\end{align}
Let $\widehat{R}^{\uu\ww}_{123}$ be the cluster transformation $\widehat{R}_{123}$ \eqref{dode}
viewed as the one for the canonical variables~${\{\uu_i, \ww_i\}_{i=1,2,3}}$.
%
Then from \eqref{Yw}, \eqref{Ypw} and \eqref{Pbch0}, we have
\begin{align}\label{knm}
\widehat{R}^{\uu\ww}_{123} &= \mathrm{Ad}\bigl(R(\mathscr{P}_1,\mathscr{P}_2,\mathscr{P}_3)_{123}\bigr).
\end{align}
Note that the right-hand side is invariant under $R \rightarrow cR$ for any scalar $c$.
A proper normalization of $R$ based on a symmetry consideration will be proposed
in Section \ref{ss:symR}.

Formally the results \eqref{uwL1} and \eqref{knm} may be stated as the commutativity of the diagrams
\begin{align*}%\label{Rcom}
\begin{CD}
\mathcal{Y}\bigl(B'\bigr) @> {\tau_{--++}}>> \mathcal{Y}(B) \\
@V{\phi'_\mathrm{SB}}VV @VV{\phi_\mathrm{SB}}V \\
\mathcal{A}_3 @>{\tau^{\uu\ww}_{--++}}>> \mathcal{A}_3,
\end{CD}
\qquad\qquad
\begin{CD}
\mathcal{Y}\bigl(B'\bigr) @> {\widehat{R}_{123}}>> \mathcal{Y}(B) \\
@V{\phi'_\mathrm{SB}}VV @VV{\phi_\mathrm{SB}}V \\
\mathcal{A}_3 @>{\widehat{R}^{\uu\ww}_{123}}>> \mathcal{A}_3.
\end{CD}
\end{align*}

Extending \eqref{RL0} and \eqref{RL1}, we introduce
$R_{ijk}=R(\mathscr{P}_i ,\mathscr{P}_j ,\mathscr{P}_k)_{ijk}$
by
\begin{align}
R_{ijk}
={}&
\Psi_q\bigl(\e^{-d_i-c_j-b_k+\uu_i+\uu_k+\ww_i-\ww_j+\ww_k}\bigr)^{-1}
\Psi_q\bigl(\e^{-d_i-c_j-b_k-e_k+\uu_i-\uu_k+\ww_i-\ww_j+\ww_k}\bigr)^{-1}
\nonumber \\
& \times
\Psi_q\bigl(\e^{d_i+e_i+c_j+b_k+\uu_i-\uu_k-\ww_i+\ww_j-\ww_k}\bigr)\nonumber \\
& \times
\Psi_q\bigl(\e^{d_i+e_i+c_j+e_j+b_k+\uu_i+2\uu_j-\uu_k-\ww_i+\ww_j-\ww_k}\bigr)P_{ijk}
\label{rijk0}\\
={}&
\Psi_q\bigl(\e^{-d_i-c_j-b_k+\uu_i+\uu_k+\ww_i-\ww_j+\ww_k}\bigr)^{-1}
\Psi_q\bigl(\e^{-d_i-c_j-b_k-e_k+\uu_i-\uu_k+\ww_i-\ww_j+\ww_k}\bigr)^{-1}
\nonumber\\
& \times P_{ijk}
\Psi_q\bigl(\e^{-b_i-a_j-d_k-e_k+\uu_i-\uu_k+\ww_i-\ww_j+\ww_k}\bigr)
\Psi_q\bigl(\e^{-b_i-a_j-d_k+\uu_i+\uu_k+\ww_i-\ww_j+\ww_k}\bigr),
\label{rijk}
\end{align}
where $P_{ijk}$ is given by \eqref{pijk}.
Now we state the main result of the paper.

\begin{Theorem}\label{th:main}
$R(\mathscr{P}_i, \mathscr{P}_j, \mathscr{P}_k)_{ijk} $ satisfies the tetrahedron equation
\begin{gather}
 R(\mathscr{P}_1, \mathscr{P}_2, \mathscr{P}_4)_{124}
R(\mathscr{P}_1, \mathscr{P}_3, \mathscr{P}_5)_{135}
R(\mathscr{P}_2, \mathscr{P}_3, \mathscr{P}_6)_{236}
R(\mathscr{P}_4, \mathscr{P}_5, \mathscr{P}_6)_{456}\nonumber
\\
\qquad= R(\mathscr{P}_4, \mathscr{P}_5, \mathscr{P}_6)_{456}
R(\mathscr{P}_2, \mathscr{P}_3, \mathscr{P}_6)_{236}
R(\mathscr{P}_1, \mathscr{P}_3, \mathscr{P}_5)_{135}
R(\mathscr{P}_1, \mathscr{P}_2, \mathscr{P}_4)_{124}.\label{tela}
\end{gather}
\end{Theorem}
\begin{proof}
 Consider the dilogarithm identity \eqref{pseq} with
 $(\varepsilon_1,\varepsilon_2,\varepsilon_3,\varepsilon_4) =(-, -,
 +, +)$ in terms of canonical variables\footnote{The signs
 $\varepsilon_1$, $\varepsilon_2$, $\varepsilon_3$, $\varepsilon_4$
 are actually $-$, $-$, $+$, $+$, but for clarity, they are left as
 they are.}
\begin{align}\label{zt16}
\Psi_q\bigl(\tilde{Z}_1\bigr)^{\varepsilon_1}\dotsm \Psi_q\bigl({\tilde Z}_{16}\bigr)^{\varepsilon_4}
=
\Psi_q\bigl(\tilde{Z'}_1\bigr)^{\varepsilon_1}\dotsm \Psi_q\bigl(\tilde{Z'}_{16}\bigr)^{\varepsilon_4}.
\end{align}
Here $\tilde{Z}_i = \tilde{\phi}_\mathrm{SB}(Z_i)$ and
$\tilde{Z'}_i =\tilde{\phi}'_\mathrm{SB}(Z'_i)$ with $Z_i$ and $Z'_i$
given in \eqref{zdL} and \eqref{zdR}. The morphisms~$\tilde{\phi}_\mathrm{SB}$ and $\tilde{\phi}'_\mathrm{SB}$ are natural
generalizations of \eqref{Yw} and \eqref{Ypw}. They send the
$Y$-variables to $\mathcal{W}_6$ according to the rule in Figure
\ref{fig:para} applied to the bottom right diagram in Figure
\ref{fig:bTE}.\footnote{The map $\phi_\mathrm{SB}$ does not spoil the
 well-definedness of the expansion like \eqref{c16} with respect to
 $\e^{\uu_i}$ and $\e^{\ww_i}$ since it preserves the rank of the quantum
 torus generated by $Y_i$ ($i=3, 6 , 7 , 8 , 11 , 12 , 13 ,
 14 , 15$).}
%
Multiplication of \eqref{pte} to \eqref{zt16} from the right leads to
\begin{gather}
\Psi_q\bigl(\tilde{Z}_1\bigr)^{\varepsilon_1}\dotsm \Psi_q\bigl({\tilde Z}_{16}\bigr)^{\varepsilon_4}
P_{124}P_{135}P_{236}P_{456}\nonumber\\
\qquad=
\Psi_q\bigl(\tilde{Z'}_1\bigr)^{\varepsilon_1}\dotsm \Psi_q\bigl(\tilde{Z'}_{16}\bigr)^{\varepsilon_4}
P_{456}P_{236}P_{135}P_{124}.\label{zp16}
\end{gather}
Let us consider the left-hand side.
It is obviously equal to
\begin{gather}
\Psi_q\bigl(\tilde{Z}_1\bigr)^{\varepsilon_1}\dotsm \Psi_q\bigl({\tilde Z}_{4}\bigr)^{\varepsilon_4}
P_{124}
P^{-1}_{124}
\Psi_q\bigl(\tilde{Z}_5\bigr)^{\varepsilon_1}\dotsm \Psi_q\bigl({\tilde Z}_{8}\bigr)^{\varepsilon_4}
P_{124} \cdot P_{135}\nonumber
\\
\qquad \times
P^{-1}_{135} P^{-1}_{124}
\Psi_q\bigl(\tilde{Z}_9\bigr)^{\varepsilon_1}\dotsm \Psi_q\bigl({\tilde Z}_{12}\bigr)^{\varepsilon_4}
P_{124} P_{135} \cdot P_{236}\nonumber
\\
\qquad \times
P^{-1}_{236} P^{-1}_{135} P^{-1}_{124}
\Psi_q\bigl(\tilde{Z}_{13}\bigr)^{\varepsilon_1}\dotsm \Psi_q\bigl({\tilde Z}_{16}\bigr)^{\varepsilon_4}
P_{124} P_{135} P_{236} \cdot P_{456}.\label{tochu1}
\end{gather}
On the other hand from \eqref{tijk} and the image of $\eqref{zu}$ by $\tilde{\phi}_\mathrm{SB}$, we know
\begin{align*}
{\tilde U}_i = \begin{cases}
\tilde{Z}_i, & i=1,\dotsc, 4,
\\
P_{124}^{-1} \tilde{Z}_i P_{124}, & i=5,\dotsc, 8,
\\
P_{135}^{-1}P_{124}^{-1} \tilde{Z}_i P_{124}P_{135}, & i=9,\dotsc, 12,
\\
P_{236}^{-1}P_{135}^{-1}P_{124}^{-1} \tilde{Z}_i
P_{124}P_{135} P_{236}, & i=13,\dotsc, 16,
\end{cases}
\end{align*}
where $\tilde{U}_i = \tilde{\phi}_\mathrm{SB}(U_i)$.
Thus \eqref{tochu1} is cast into the form
\begin{gather*}
\Psi_q\bigl(\tilde{U}_1\bigr)^{\varepsilon_1}\cdots \Psi_q\bigl({\tilde U}_{4}\bigr)^{\varepsilon_4}
P_{124}
\Psi_q\bigl(\tilde{U}_5\bigr)^{\varepsilon_1}\cdots \Psi_q\bigl({\tilde U}_{8}\bigr)^{\varepsilon_4}
P_{135}
\\
\qquad\times
\Psi_q\bigl(\tilde{U}_9\bigr)^{\varepsilon_1}\cdots \Psi_q\bigl({\tilde U}_{12}\bigr)^{\varepsilon_4}
P_{236}
\Psi_q\bigl(\tilde{U}_{13}\bigr)^{\varepsilon_1}\cdots \Psi_q\bigl({\tilde U}_{16}\bigr)^{\varepsilon_4}
P_{456}.
\end{gather*}
This is identified with the left-hand side of \eqref{tela} for \eqref{rijk0}.
The right-hand side of \eqref{tela} is similarly derived from that in \eqref{zp16}.
\end{proof}
